\chapter{Lebesgue可測関数}

\paragraph{本章の目標}
ここまでで，どの集合に面積を与えられるかが分かった．
次は，その上でどの関数を積分できるかを決める必要がある．
そのための基本概念がLebesgue可測関数である．
本章では
\begin{enumerate}
    \item Lebesgue可測関数の定義
    \item 可測写像の定義が逆像で書かれる理由
    \item 連続関数・指示関数などの典型例
    \item 和・積・逆数・合成・sup/inf/limsup/liminf/lim などの可測性
    \item almost everywhere による同値関係
\end{enumerate}
を証明する．

\section{動機づけ}

積分をしたいのは関数である．
しかしLebesgue積分では，まず集合に測度を与え，そのあとで関数を積分する．
したがって関数についても，
\[
\{x \mid f(x)>a\},\qquad
\{x \mid f(x)\ge a\}
\]
のようなレベル集合がちゃんと測れることが必要になる．

Riemann積分では，定義域を細かく分割して関数値を調べた．
それに対してLebesgue積分では，
まず関数値 $a$ を固定し，
\[
f(x)>a
\]
となる点の集まりを調べる．
つまり，関数そのものではなく，
関数がどの集合の上でどのような値をとるかを
測度論の言葉で記述するのである．

\section{可測写像}

\begin{definition}[可測空間]
集合 $X$ とその部分集合族 $\mathcal M$ の組
\[
(X,\mathcal M)
\]
で，$\mathcal M$ が $X$ 上の $\sigma$-加法族であるものを可測空間という．
\end{definition}

\begin{definition}[可測写像]
可測空間 $(X,\mathcal M)$ と $(Y,\mathcal N)$ の間の写像
\[
f:X\to Y
\]
が可測であるとは，任意の $B \in \mathcal N$ に対して
\[
f^{-1}(B)\in \mathcal M
\]
が成り立つことをいう．
\end{definition}

\begin{proposition}[逆像の基本性質]
任意の写像 $f:X\to Y$ と部分集合族 $\{B_\alpha\}_{\alpha\in\Gamma}\subset \mathcal P(Y)$ に対して
次が成り立つ．
\begin{enumerate}
    \item
    \[
    f^{-1}(Y\setminus B)=X\setminus f^{-1}(B)
    \]
    \item
    \[
    f^{-1}\left(\bigcup_{\alpha\in\Gamma}B_\alpha\right)
    =
    \bigcup_{\alpha\in\Gamma}f^{-1}(B_\alpha)
    \]
    \item
    \[
    f^{-1}\left(\bigcap_{\alpha\in\Gamma}B_\alpha\right)
    =
    \bigcap_{\alpha\in\Gamma}f^{-1}(B_\alpha)
    \]
\end{enumerate}
\end{proposition}
\begin{proof}
\textbf{(1)}\quad
$x\in f^{-1}(Y\setminus B)$ であることは
\[
f(x)\in Y\setminus B
\]
と同値であり，これは
\[
f(x)\notin B
\]
と同値である．
さらにこれは
\[
x\notin f^{-1}(B)
\]
と同値だから
\[
x\in X\setminus f^{-1}(B)
\]
である．
したがって
\[
f^{-1}(Y\setminus B)=X\setminus f^{-1}(B).
\]

\textbf{(2)}\quad
$x\in f^{-1}(\bigcup_\alpha B_\alpha)$ であることは
\[
f(x)\in \bigcup_{\alpha\in\Gamma}B_\alpha
\]
と同値であり，これは
ある $\alpha$ が存在して
\[
f(x)\in B_\alpha
\]
となることと同値である．
さらにこれは
ある $\alpha$ が存在して
\[
x\in f^{-1}(B_\alpha)
\]
となることと同値である．
したがって
\[
x\in \bigcup_{\alpha\in\Gamma}f^{-1}(B_\alpha)
\]
である．

\textbf{(3)}\quad
(1), (2) と De Morgan の法則から従う．
\end{proof}

\begin{remark}[可測写像の定義が逆なのはなぜか]
連続写像も可測写像も，
定義はどちらも像ではなく逆像で書かれる．
その理由は，逆像が
\[
\text{補集合},\qquad \text{可算和},\qquad \text{可算共通部分}
\]
を正確に保つからである．
実際，直前の命題より
\[
f^{-1}(Y\setminus B)=X\setminus f^{-1}(B),
\qquad
f^{-1}\left(\bigcup_n B_n\right)=\bigcup_n f^{-1}(B_n)
\]
が成り立つ．
したがって，$Y$ 側の $\sigma$-加法族の構造が
$X$ 側へそのまま引き戻される．

これに対して像 $f(B)$ は，一般には
\[
f(X\setminus A)=Y\setminus f(A)
\]
も
\[
f\left(\bigcap_n A_n\right)=\bigcap_n f(A_n)
\]
も満たさない．
したがって「構造を保つ」という意味で自然なのは像ではなく逆像なのである．
\end{remark}

\section{Lebesgue可測関数の定義}

\begin{definition}[制限された $\sigma$-加法族]
$E \subset \RR^d$ をLebesgue可測集合とする．
\[
\mathcal L_E^d:=\{E\cap A \mid A\in \mathcal L^d\}
\]
を $E$ 上の制限Lebesgue $\sigma$-加法族という．
\end{definition}

\begin{definition}[Borel集合族]
\[
\mathcal B(\eRR)
\]
を拡大実数直線 $\eRR$ のBorel集合族とする．
すなわち，開集合全体が生成する $\sigma$-加法族である．
\end{definition}

\begin{definition}[Lebesgue可測関数]
$E \subset \RR^d$ をLebesgue可測集合とする．
写像
\[
f:E\to \eRR
\]
がLebesgue可測であるとは，
\[
f:(E,\mathcal L_E^d)\to (\eRR,\mathcal B(\eRR))
\]
が可測写像であることをいう．
\end{definition}

\begin{theorem}[判定法]
$f:E\to\eRR$ に対して，次は同値である．
\begin{enumerate}
    \item $f$ はLebesgue可測である
    \item 任意の $a\in\RR$ に対して
    \[
    \{x\in E \mid f(x)>a\}
    \]
    は可測である
    \item 任意の $a\in\RR$ に対して
    \[
    \{x\in E \mid f(x)\ge a\}
    \]
    は可測である
    \item 任意の $a\in\RR$ に対して
    \[
    \{x\in E \mid f(x)<a\}
    \]
    は可測である
    \item 任意の $a\in\RR$ に対して
    \[
    \{x\in E \mid f(x)\le a\}
    \]
    は可測である
\end{enumerate}
\end{theorem}
\begin{proof}
\textbf{(1) $\Rightarrow$ (2)}\quad
\[
(a,\infty]\in \mathcal B(\eRR)
\]
だから，可測性の定義より
\[
f^{-1}((a,\infty])
=
\{x\in E \mid f(x)>a\}
\]
は可測である．

\textbf{(2) $\Rightarrow$ (3)}\quad
任意の $a\in\RR$ に対して
\[
\{f\ge a\}
=
\bigcap_{n=1}^{\infty}\left\{f>a-\frac1n\right\}
\]
である．
右辺は可測集合の可算共通部分だから可測である．

\textbf{(3) $\Rightarrow$ (4)}\quad
\[
\{f<a\}=E\setminus\{f\ge a\}
\]
だから可測である．

\textbf{(4) $\Rightarrow$ (5)}\quad
\[
\{f\le a\}
=
\bigcap_{n=1}^{\infty}\left\{f<a+\frac1n\right\}
\]
だから可測である．

\textbf{(5) $\Rightarrow$ (1)}\quad
\[
\mathcal C:=\{B\subset \eRR \mid f^{-1}(B)\in \mathcal L_E^d\}
\]
とおく．
逆像の基本性質より，$\mathcal C$ は $\eRR$ 上の $\sigma$-加法族である．
仮定より
\[
(-\infty,a]\in \mathcal C
\qquad (a\in\RR)
\]
である．
ところが
\[
\mathcal B(\eRR)=\sigma(\{(-\infty,a]\mid a\in\RR\})
\]
だから
\[
\mathcal B(\eRR)\subset \mathcal C
\]
である．
したがって $f$ は可測である．
\end{proof}

\section{基本例}

\begin{example}[連続関数]
$E \subset \RR^d$ 上の連続関数
\[
f:E\to\RR
\]
はLebesgue可測である．
\end{example}
\begin{proof}
任意の $a\in\RR$ に対して
\[
\{x\in E \mid f(x)>a\}=E\cap f^{-1}((a,\infty))
\]
である．
$f$ が連続だから $f^{-1}((a,\infty))$ は $\RR^d$ の開集合である．
したがって右辺は $E$ の可測部分集合である．
よって判定法より $f$ は可測である．
\end{proof}

\begin{example}[指示関数]
$A \subset E$ が可測なら
\[
1_A:E\to\RR
\]
はLebesgue可測である．
\end{example}
\begin{proof}
任意の $a\in\RR$ に対して
\[
\{1_A>a\}
=
\begin{cases}
\emptyset & (a\ge 1),\\
A & (0\le a<1),\\
E & (a<0).
\end{cases}
\]
いずれも可測だから，判定法より $1_A$ は可測である．
\end{proof}

\begin{example}[Dirichlet関数]
\[
f=1_{\QQ\cap[0,1]}
\]
は $[0,1]$ 上Lebesgue可測である．
\end{example}
\begin{proof}
第5章で
\[
\QQ\cap[0,1]
\]
はLebesgue可測であることを示した．
したがって直前の例より $f$ は可測である．
\end{proof}

\begin{remark}
Dirichlet関数はRiemann積分できなかったが，
Lebesgue積分論ではまず「可測である」ことが分かる．
積分できるかどうかは次章で考える．
\end{remark}

\section{ベクトル値可測写像と合成}

\begin{proposition}[有限個の可測関数をまとめた写像]
$f_1,\dots,f_n:E\to\RR$ を可測関数とする．
このとき
\[
F:E\to\RR^n,
\qquad
F(x):=(f_1(x),\dots,f_n(x))
\]
はBorel可測である．
\end{proposition}
\begin{proof}
\[
\mathcal C:=\{B\in\mathcal B(\RR^n)\mid F^{-1}(B)\in \mathcal L_E^d\}
\]
とおく．
逆像の基本性質より $\mathcal C$ は $\sigma$-加法族である．

任意の $a_1,\dots,a_n\in\RR$ に対して
\begin{align*}
F^{-1}\bigl((-\infty,a_1]\times\cdots\times(-\infty,a_n]\bigr)
=
\bigcap_{k=1}^n \{x\in E\mid f_k(x)\le a_k\}
\end{align*}
であり，右辺は可測である．
したがって
\[
(-\infty,a_1]\times\cdots\times(-\infty,a_n]\in\mathcal C.
\]
ところが，このような直方体は $\mathcal B(\RR^n)$ を生成するから
\[
\mathcal B(\RR^n)\subset \mathcal C
\]
である．
よって $F$ は可測である．
\end{proof}

\begin{theorem}[連続写像との合成]
$F:E\to\RR^n$ を可測写像，
\[
\Phi:\RR^n\to\RR
\]
を連続関数とする．
このとき
\[
\Phi\circ F:E\to\RR
\]
は可測である．
\end{theorem}
\begin{proof}
任意のBorel集合 $B\subset\RR$ に対して，
$\Phi$ の連続性より
\[
\Phi^{-1}(B)\in\mathcal B(\RR^n)
\]
である．
したがって
\[
(\Phi\circ F)^{-1}(B)=F^{-1}(\Phi^{-1}(B))
\]
は可測である．
よって $\Phi\circ F$ は可測である．
\end{proof}

\begin{corollary}[代数演算と合成]
$f,g:E\to\RR$ を可測関数とする．
このとき次が成り立つ．
\begin{enumerate}
    \item 任意の $c\in\RR$ に対して $cf$ は可測である
    \item $f+g$ と $fg$ は可測である
    \item $\max(f,g)$ と $\min(f,g)$ は可測である
    \item $p>0$ に対して $|f|^p$ は可測である
    \item $\phi:\RR\to\RR$ が連続なら $\phi\circ f$ は可測である
\end{enumerate}
\end{corollary}
\begin{proof}
\begin{enumerate}
    \item
    \[
    \Phi(t):=ct
    \]
    は連続だから，直前の定理を $n=1$ に適用すればよい．

    \item
    \[
    \Phi(u,v):=u+v,\qquad \Psi(u,v):=uv
    \]
    は $\RR^2$ 上連続である．
    また $(f,g):E\to\RR^2$ は前命題より可測である．
    よって
    \[
    f+g=\Phi\circ(f,g),\qquad fg=\Psi\circ(f,g)
    \]
    は可測である．

    \item
    \[
    \Phi(u,v):=\max(u,v),\qquad \Psi(u,v):=\min(u,v)
    \]
    は連続だから，同様に従う．

    \item
    \[
    \phi(t):=|t|^p
    \]
    は連続だから従う．

    \item 直前の定理そのものである．
\end{enumerate}
\end{proof}

\begin{proposition}[逆数]
$f:E\to\RR$ を可測関数とする．
このとき
\[
E_0:=\{x\in E\mid f(x)\ne 0\}
\]
は可測であり，
\[
\frac{1}{f}:E_0\to\RR
\]
は可測である．
特に $f(x)\ne 0$ がすべての $x\in E$ で成り立つなら，
\[
1/f
\]
は $E$ 上可測である．
\end{proposition}
\begin{proof}
まず
\[
E_0=E\setminus\{f=0\}
\]
だから可測である．

次に
\[
\phi:\RR\setminus\{0\}\to\RR,
\qquad
\phi(t)=\frac1t
\]
は連続である．
$f$ を $E_0$ に制限した写像
\[
f|_{E_0}:E_0\to \RR\setminus\{0\}
\]
は可測であるから，連続写像との合成より
\[
\phi\circ f|_{E_0}=\frac1f
\]
は可測である．
\end{proof}

\section{不等式・極限で作られる関数}

\begin{proposition}[比較からできる集合]
$f,g:E\to\eRR$ を可測関数とする．
このとき次の集合はすべて可測である．
\begin{enumerate}
    \item
    \[
    \{x\in E \mid f(x)>g(x)\}
    \]
    \item
    \[
    \{x\in E \mid f(x)\ge g(x)\}
    \]
    \item
    \[
    \{x\in E \mid f(x)=g(x)\}
    \]
\end{enumerate}
\end{proposition}
\begin{proof}
\textbf{(1)}\quad
まず
\[
\{f>g\}
=
\bigcup_{q\in\QQ}\bigl(\{f>q\}\cap\{g<q\}\bigr)
\]
を示す．
左辺に $x$ が属するとする．
すると $f(x)>g(x)$ だから，有理数の稠密性より
\[
g(x)<q<f(x)
\]
を満たす $q\in\QQ$ が存在する．
したがって $x$ は右辺に属する．
逆向きは明らかである．
右辺は可測集合の可算和だから可測である．

\textbf{(2)}\quad
\[
\{f\ge g\}=E\setminus\{g>f\}
\]
だから可測である．

\textbf{(3)}\quad
\[
\{f=g\}=\{f\ge g\}\cap\{g\ge f\}
\]
だから可測である．
\end{proof}

\begin{proposition}[可算個の上限と下限]
$f_n:E\to\eRR$ を可測関数列とする．
このとき
\[
\sup_{n}f_n,\qquad \inf_{n}f_n
\]
は可測である．
\end{proposition}
\begin{proof}
任意の $a\in\RR$ に対して
\[
\left\{\sup_n f_n>a\right\}
=
\bigcup_{n=1}^{\infty}\{f_n>a\}
\]
である．
右辺は可測集合の可算和だから可測である．
よって $\sup_n f_n$ は可測である．

また
\[
\inf_n f_n=-\sup_n(-f_n)
\]
であり，各 $-f_n$ も可測だから，$\inf_n f_n$ も可測である．
\end{proof}

\begin{proposition}[limsup, liminf, limit]
$f_n:E\to\eRR$ を可測関数列とする．
このとき
\[
\limsup_{n\to\infty}f_n,\qquad
\liminf_{n\to\infty}f_n
\]
は可測である．
さらに，もし各点で極限
\[
f(x)=\lim_{n\to\infty}f_n(x)
\]
が存在するなら，$f$ も可測である．
\end{proposition}
\begin{proof}
定義により
\[
\limsup_{n\to\infty}f_n
=
\inf_{n\in\NN}\sup_{k\ge n}f_k,
\qquad
\liminf_{n\to\infty}f_n
=
\sup_{n\in\NN}\inf_{k\ge n}f_k.
\]
右辺は可測関数に対する可算上限・可算下限で作られているから可測である．

さらに各点で極限が存在するなら
\[
f=\limsup_{n\to\infty}f_n=\liminf_{n\to\infty}f_n
\]
だから $f$ は可測である．
\end{proof}

\section{almost everywhere}

\begin{definition}[almost everywhere]
$E \subset \RR^d$ を可測集合とし，$f,g:E\to\eRR$ を関数とする．
\[
f=g \quad \text{a.e.\ on } E
\]
とは
\[
\{x\in E\mid f(x)\ne g(x)\}
\]
がLebesgue零集合であることをいう．
\end{definition}

\begin{proposition}
$E$ 上の関数全体において，
\[
f\sim g \iff f=g \text{ a.e.\ on } E
\]
で定める関係 $\sim$ は同値関係である．
\end{proposition}
\begin{proof}
\textbf{反射律}\quad
\[
\{f\ne f\}=\emptyset
\]
だから $f\sim f$ である．

\textbf{対称律}\quad
\[
\{f\ne g\}=\{g\ne f\}
\]
だから $f\sim g$ なら $g\sim f$ である．

\textbf{推移律}\quad
$f\sim g$ かつ $g\sim h$ とする．
すると
\[
\{f\ne h\}\subset \{f\ne g\}\cup\{g\ne h\}
\]
である．
右辺は零集合の可算和だから零集合である．
したがって $f\sim h$ である．
\end{proof}

\begin{proposition}[a.e.\ 修正による可測性]
$f:E\to\eRR$ を可測関数とし，
$g:E\to\eRR$ が
\[
f=g \quad \text{a.e.\ on } E
\]
を満たすとする．
このとき $g$ も可測である．
\end{proposition}
\begin{proof}
\[
N:=\{x\in E\mid f(x)\ne g(x)\}
\]
とおくと，$N$ は零集合である．
Lebesgue測度は完備だから，$N$ の任意の部分集合は可測である．

任意の $a\in\RR$ に対して
\[
\{g>a\}
=
\bigl(\{f>a\}\cap(E\setminus N)\bigr)\cup\bigl(\{g>a\}\cap N\bigr)
\]
である．
右辺第1項は可測であり，
第2項は $N$ の部分集合だから可測である．
よって
\[
\{g>a\}
\]
は可測である．
したがって判定法より $g$ は可測である．
\end{proof}

\begin{corollary}[a.e.\ 極限の可測性]
$f_n:E\to\eRR$ を可測関数列とし，
\[
f_n(x)\to f(x)\quad \text{a.e.\ on } E
\]
とする．
このとき $f$ は可測である．
\end{corollary}
\begin{proof}
\[
g:=\limsup_{n\to\infty}f_n
\]
とおくと，前節より $g$ は可測である．
また極限が存在する点では
\[
g(x)=\lim_{n\to\infty}f_n(x)=f(x)
\]
である．
したがって
\[
f=g \quad \text{a.e.\ on } E.
\]
ゆえに直前の命題より $f$ は可測である．
\end{proof}

\section{解釈とまとめ}

Lebesgue可測関数とは，
\[
\{f>a\}
\]
のようなレベル集合が可測である関数であった．
つまり，関数の値そのものより，
その値をとる領域の幾何学が測度で追えることが本質である．

本章で得た重要な事実は次の通りである．
\begin{enumerate}
    \item 連続関数と指示関数は可測である
    \item 可測関数は和・積・逆数・合成で閉じている
    \item $\sup,\inf,\limsup,\liminf,\lim$ をとっても可測性は保たれる
    \item 零集合上で値を変えても可測性は変わらない
\end{enumerate}

次章では，この可測関数に対してLebesgue積分を定義する．
その出発点は，
\[
\text{指示関数} \to \text{単関数} \to \text{一般の可測関数}
\]
という拡張である．
